\documentclass[sigplan,nonacm]{acmart}
\makeatletter
\renewcommand\@formatdoi[1]{\ignorespaces}
\makeatother
\settopmatter{printfolios=false,printccs=false,printacmref=false}

\setcopyright{none}
\acmConference[]{}{}{}

\usepackage{enumitem}
\usepackage{mathtools}
\usepackage{amsmath}

\usepackage{listings}

\lstdefinelanguage{pcpp}{%
  language     = C++,
  morekeywords = {type,then,of,match,with,null},
  keywordstyle=\ttfamily\bfseries
}

\usepackage{natbib}
\usepackage{url}

\newcommand{\defeq}{\stackrel{\text{def}}{=}}
\newcommand{\setbuild}[2]{\{#1~|~#2\}}

% Syntax highlighting
\newdimen\zzlistingsize
\newdimen\zzlistingsizedefault
\zzlistingsizedefault=9pt
\newdimen\kwlistingsize
\kwlistingsize=9pt
\zzlistingsize=\zzlistingsizedefault
\gdef\lco{black}
\newcommand{\lcRed}{\gdef\lco{nred}}
\newcommand{\lcBlack}{\gdef\lco{black}}
\newcommand{\lcNormal}{\lcBlack}
\newcommand{\lcBlue}{\gdef\lco{nblue}}
\newcommand{\lcGreen}{\gdef\lco{ngreen}}
\newcommand{\lcPurple}{\gdef\lco{npurple}}
\newcommand{\Lstbasicstyle}{\fontsize{\zzlistingsize}{1.1\zzlistingsize}\ttfamily\color{\lco}}
\newcommand{\keywordstyle}{\fontsize{1.09\kwlistingsize}{\kwlistingsize}\normalfont\bf\color{\lco}}
\newlength{\zzlstwidth}
\settowidth{\zzlstwidth}{{\Lstbasicstyle~}}
\newcommand{\lcm}{\color{\lco}}

\lstset{
  basewidth=\zzlstwidth,
%  basicstyle=\Lstbasicstyle,
  basicstyle=\small\ttfamily,
%  identifierstyle=,
  numbers=left, numberstyle=\tiny, numbersep=5pt,
%  frame=single,
  keywordstyle=\keywordstyle,
  keywords={fun,cfun,sfun,read,case,of,type,functor,structure,signature,sig,struct,datatype,let,letpar,val,fn,in,end,as,select,mod,apply,sapply,capply,write},  %nil,cons,
  commentstyle=\rmfamily\slshape\color{\lco},
  morecomment=[n]{(*}{*)},
  lineskip={-0pt},
%  columns=[l]flexible,
  columns=[l]fullflexible,
  keepspaces=true,
  mathescape=true,
% NOTE: need TWO sets of braces around each definition below!
  literate={->}{{$\lcm\rightarrow$}}2
           {->^S}{{$\lcm\hspace{-1pt}\arrowsymbol{\S}\hspace{-1pt}$}}2
           {->^C}{{$\lcm\hspace{-1pt}\arrowsymbol{\C}\hspace{-1pt}$}}2
           {->^D}{{$\lcm\hspace{-1pt}\arrowsymbol{\delta}\hspace{-1pt}$}}2
           {->^Dh}{{$\lcm\hspace{-1pt}\arrowsymbol{\delta_h}\hspace{-1pt}$}}2
           {->^Dt}{{$\lcm\hspace{-1pt}\arrowsymbol{\delta_t}\hspace{-1pt}$}}2
%           {-c->}{${c\rightarrow}$}3
%           {-s->}{${s\rightarrow}$}3
%           {-c->}{$\rightarrow$}1
%           {-s->}{$\rightarrow$}1
%           {_}{{\lcm\_}}1
           {'a}{{$\alpha$}}1
           {'b}{{$\beta$}}1
           {'d}{{$\lcm\delta$}}1
}


\lstset{language=pcpp,numbers=none}

\newcommand{\vz}{\mathtt{v_0}}
\newcommand{\vo}{\mathtt{v_1}}
\newcommand{\vt}{\mathtt{v_2}}
\newcommand{\fo}{\mathtt{f_1}}
\newcommand{\ft}{\mathtt{f_2}}
% A step number, flush right in a listing.
\newcommand{\stepc}[1]{\hfill{\rmfamily\slshape // (#1)}}
% A timeline: two columns of pseudocode, in the font of listings.
\newenvironment{timeline}[1]
  {\par\raggedright\noindent\small\ttfamily\begin{tabular}{@{}#1@{}}}
  {\end{tabular}\par}
\newcommand{\tlhead}[1]{{\rmfamily\itshape #1}}
\newcommand{\ind}{\hphantom{(0)~}}

\title{Native fork join: work stealing by context switching}
% TODO: add co-author, and check affiliation.
\author{Mike Rainey}
\affiliation{
  \institution{Carnegie Mellon University}
  \city{Pittsburgh}
  \state{PA}
  \country{USA}
}
\email{me@mike-rainey.site}

\begin{document}

\begin{abstract}
TaskPaRTS is a C++ library for fork-join parallelism that requires no
compiler support.
%
At its core is a technique we call \emph{native fork join}, which
delivers work-stealing scheduling by context switching.
%
A fork saves the continuation of the parent task as a handful of
machine registers that point into the parent's own stack.
%
When the join finds that the second branch was stolen, the worker
neither blocks nor runs other work on top of the parent's frame:
it leaves the parent's stack suspended and returns to its scheduler.
%
Whichever worker finishes the last branch then resumes the parent,
possibly on a different thread.
%
When nothing is stolen, both branches run on the parent's stack, like
ordinary calls, and the fork allocates no memory.
%
In this short paper, we present the algorithm in pseudocode, explain
why it is safe, and describe what it means for programs that use it.
%
A companion implementation, micro taskparts, implements the same
algorithm in a few hundred lines of C++.
\end{abstract}

\maketitle

\section{Challenge: the join, in a library}

We are to implement binary fork join, which runs two thunks, possibly
in parallel, and returns once both have finished.
%
\begin{lstlisting}
fork2join($\fo$, $\ft$ : () -> void) -> void
\end{lstlisting}
%
For example:
%
\begin{lstlisting}
fib(n : int) -> int {
  if (n < 2) return n
  fork2join($\lambda$ () $\Rightarrow$ {a = fib(n-1)},
            $\lambda$ () $\Rightarrow$ {b = fib(n-2)})
  return a + b }
\end{lstlisting}
%
Work stealing~\cite{halstead84,blumofele99} is well suited to
scheduling such programs.
%
Our constraint is that we are writing a library: there is no compiler
support, and client code is ordinary C++, with lambdas for thunks.

A library work stealer typically implements \lstinline|fork2join| by
\emph{child stealing}: the second branch becomes a task that thieves
may take.
%
\begin{lstlisting}
fork2join($\fo$, $\ft$ : () -> void) -> void {
  t = new_task($\ft$); push(t)  // t is now stealable
  $\fo$()
  if (pop() == t) { $\ft$(); return }  // not stolen
  while (!t.done)            // stolen: wait for thief
    { u = steal(); if (u != null) u.run() } }
\end{lstlisting}
%
The common case, where \lstinline|t| is not stolen, is cheap.
%
The difficulty is the join when \lstinline|t| was stolen.
%
The parent's frame sits on the worker's thread stack, and must stay
there until the join completes.
%
So the worker can either block, and sit idle, or, as TBB and
ParlayLib do~\cite{threadingbuildingblocksmanual,blelloch20parlaylib},
\emph{help}: run other tasks on top of the blocked frame, as in the
loop above.
%
Helping keeps the worker busy, but the parent can resume only once the
worker has unwound back down to its frame, even if the thief finished
long before.
%
Cilk~\cite{frigo98cilk5,schardl2017tapir} avoids waiting by stealing
the parent's continuation instead of the child.
%
But then the thief must resume a frame that sits on the victim's
stack, which in Cilk relies on compiler support and a cactus
stack~\cite{lee2010cilk-m}.

Native fork join, which is at the core of
TaskPaRTS~\cite{taskparts}, takes a different route.
%
When a worker finds at a join that the second branch was stolen, it
\emph{suspends the parent's stack} and returns to its scheduler, which
looks for other work.
%
The parent's continuation is a saved register context that points
into the suspended stack.
%
Whichever worker finishes the last branch restores this context, and
so the parent continues on that worker's thread.
%
No worker waits at a join, and no two workers ever run on the same
stack at once.

We present the algorithm in Section~\ref{sec:nfj}, trace its two cases
in Section~\ref{sec:cases}, explain why it is safe in
Section~\ref{sec:safety}, and discuss its costs and its consequences
for client code in Section~\ref{sec:costs}.
%
The appendices give machine-level details, the work-stealing deque,
and a map from the pseudocode to TaskPaRTS and to micro taskparts, a
minimal C++ implementation that accompanies this paper.

\section{Native fork join}
\label{sec:nfj}

\paragraph{Contexts}
%
Our one low-level primitive is the \emph{context}.
%
A context holds what the C calling convention requires a function
call to preserve: the callee-saved registers, the stack pointer, and a
return address.
%
That is 8 words on x86-64 and 21 on AArch64
(Appendix~\ref{sec:contexts}).
%
\lstinline|save($k$)| stores its caller's context in $k$ and returns
\lstinline|false|; \lstinline|restore($k$)| resumes that call to
\lstinline|save|, which this time returns \lstinline|true|.
%
The pair behaves like \lstinline|setjmp| and \lstinline|longjmp|,
except that the thread that restores a context need not be the one
that saved it, and the context may point into any stack.
%
To the compiler, \lstinline|save| is an ordinary call (one that returns
twice, like \lstinline|setjmp|), so it spills any live caller-saved
registers before the call.
%
As a result, a save or a restore is a handful of loads and stores,
with no system call.
%
We write \lstinline|swap($k_1$, $k_2$)| for
\lstinline|if (!save($k_1$)) restore($k_2$)|, which saves the current
context in $k_1$ and switches to $k_2$.

\paragraph{Vertices and workers}
%
Following the DAG calculus~\cite{Dag_calculus_16}, the scheduler
represents a computation as a DAG whose vertices each have an
\emph{in-counter}, which counts unfinished predecessors, and an
\emph{out-set}, which, in fork join, is a single successor
(Figure~\ref{fig:scheduler}).
%
A vertex also has a context, \lstinline|k|, and an action,
\lstinline|act|, which tells the scheduler what to do with the vertex.
%
\lstinline|Init| means that the vertex has never run;
\lstinline|Resume|, that it is suspended inside \lstinline|fork2join|,
with its context saved in \lstinline|k|; and \lstinline|Finish|, that
it will have finished the next time it hands control back to the
scheduler.

Each worker runs a scheduler loop on its thread's own stack.
%
While the worker runs a vertex, \lstinline|cur|, it keeps the loop's
context in \lstinline|k|.
%
Client code runs on separate \emph{task stacks}.
%
We write \lstinline|me| for the worker that runs the caller, which we
look up again at each use, because the worker can change across a call
to \lstinline|fork2join| (Section~\ref{sec:costs}).
%
Each worker also holds its ready vertices:
\lstinline|schedule($v$)| makes $v$ ready, \lstinline|pop()| takes back
the vertex scheduled most recently, if any, and \lstinline|steal()|
takes the oldest ready vertex of another worker, chosen at random.
%
The vertex scheduled most recently can never be stolen
(Appendix~\ref{sec:deque}), and \lstinline|fork2join| relies on this.

\begin{figure}
\hrule
\begin{lstlisting}
type action = Init | Resume | Finish
type vertex = {in : atomic int, out : vertex*,
  act : action, k : context, body : () -> void}
type worker = {k : context, cur : vertex*, ...}

loop() -> void {  // on the worker's own stack
  while (!done) {
    v = pop()
    if (v == null) {
      v = steal()
      if (v != null) schedule(v)
      continue }
    me.cur = v
    if (v.act == Init) {
      v.k = new_context(new_stack(), enter)
      v.act = Finish }
    swap(me.k, v.k)  // run v until it yields
    if (v.act == Finish) notify(v) } }

enter() -> void {  // first frame on a fresh stack
  me.cur.body(); restore(me.k) }

notify(v : vertex*) -> void {
  u = v.out  // read before u.in can reach 0
  if (u == null) { done = true; return }  // root
  if (fetch_and_add(u.in, -1) == 1) schedule(u) }
\end{lstlisting}
\hrule
\caption{The scheduler}
\label{fig:scheduler}
\end{figure}

\begin{figure}
\hrule
\begin{lstlisting}[escapechar=@]
fork2join($\fo$, $\ft$ : () -> void) -> void {
  $\vz$ = me.cur
  $\vo$ = vertex{out = $\vz$, body = $\lambda$ () $\Rightarrow$ {}}
  $\vt$ = vertex{out = $\vz$, body = $\ft$}
  if (save($\vz$.k))@\stepc{1}@
    { $\vz$.act = Finish; return }
  $\vz$.act = Resume; $\vz$.in = 2@\stepc{2}@
  schedule($\vt$); schedule($\vo$)
  $\vo$.act = Resume@\stepc{3}@
  swap($\vo$.k, me.k); $\vo$.act = Finish
  $\fo$()
  if (pop() == null) restore(me.k)@\stepc{4}@
  $\vt$.act = Resume; schedule($\vt$)@\stepc{5}@
  swap($\vt$.k, me.k); $\vt$.act = Finish
  $\ft$()
  swap($\vz$.k, me.k); $\vz$.act = Finish }@\stepc{6}@
\end{lstlisting}
\hrule
\caption{Fork join}
\label{fig:fork2join}
\end{figure}

\paragraph{The scheduler}
%
The scheduler loop (Figure~\ref{fig:scheduler}) is the only place
where vertices start, resume, and finish.
%
It pops a vertex, or else steals one and makes it ready, and then makes
the vertex current.
%
If the vertex has never run, the loop gives it a fresh stack and a
context that, once restored, calls \lstinline|enter| at the top of that
stack.
%
\lstinline|enter| runs the body of the vertex, and then switches to the
scheduler of whichever worker it is on by then.
%
The loop then switches to the vertex, and regains control when the
vertex yields.
%
If the action of the vertex is then \lstinline|Finish|, the loop
notifies its successor: it decrements the successor's in-counter, and
the worker whose decrement reaches zero schedules the successor.
%
A call to \lstinline|fork2join| from outside the scheduler starts the
workers with a root vertex that makes the call, and the computation
ends when the root, which has no successor, finishes.

\paragraph{Fork join}
%
In \lstinline|fork2join| (Figure~\ref{fig:fork2join}), the parent,
$\vz$, is the current vertex.
%
The function creates two vertices, $\vo$ and $\vt$, both with $\vz$ as
their successor.
%
$\vt$ runs $\ft$, and $\vo$ has an empty body.
%
$\vo$ stands for the parent while the parent runs $\fo$ inline, so that
any \lstinline|fork2join| nested inside $\fo$ uses the context and
in-counter of $\vo$, and leaves those of $\vz$ alone.
%
The steps are as follows.
%
\begin{enumerate}[leftmargin=*,label=(\arabic*)]
\item Save the parent's context in $\vz$\lstinline|.k|.
  If $\vt$ is stolen, this is where the parent resumes.
\item Set the in-counter of $\vz$ to 2, and schedule $\vt$, then
  $\vo$, which leaves $\vt$ stealable and $\vo$ not.
\item Yield to the scheduler, saving the context in $\vo$\lstinline|.k|.
  The scheduler pops $\vo$ and resumes it at once, so execution
  continues after the \lstinline|swap|, with $\vo$ as the current vertex.
  Run $\fo$ inline.
\item Pop.
  If the pop fails, $\vt$ was stolen.
  Then abandon the current stack by switching to the scheduler without
  saving anything.
  The stack stays suspended at step 1, and the scheduler notifies $\vz$
  that $\vo$ has finished.
\item Otherwise, the pop returned $\vt$.
  Schedule $\vt$ again, so that it is the vertex scheduled most
  recently and cannot be stolen, and yield.
  The scheduler notifies $\vz$ that $\vo$ has finished, then pops $\vt$
  and resumes it here.
  Run $\ft$ inline.
\item Yield once more, saving the context in $\vz$\lstinline|.k|.
  The scheduler notifies $\vz$ that $\vt$ has finished, which makes $\vz$
  ready, then pops $\vz$ and resumes it here, and
  \lstinline|fork2join| returns.
\end{enumerate}

\section{The two cases}
\label{sec:cases}

\begin{figure}
\hrule
\vspace{4pt}
\begin{timeline}{l@{\hspace{0.5em}}c@{\hspace{0.5em}}l}
\tlhead{on task stack S} & & \tlhead{in the scheduler of worker A} \\
(1) save($\vz$.k) \\
(2) schedule($\vt$) \\
\ind schedule($\vo$) \\
(3) swap($\vo$.k, A.k) & $\rightarrow$ & pop $\vo$, resume it \\
\ind $\fo$() \\
(4) pop() returns $\vt$ \\
(5) schedule($\vt$) \\
\ind swap($\vt$.k, A.k) & $\rightarrow$ & notify: $\vz$.in = 1 \\
& & pop $\vt$, resume it \\
\ind $\ft$() \\
(6) swap($\vz$.k, A.k) & $\rightarrow$ & notify: $\vz$.in = 0 \\
& & schedule, pop, resume $\vz$ \\
\ind \textbf{return}
\end{timeline}
\vspace{4pt}
\hrule
\vspace{4pt}
\begin{timeline}{l@{\hspace{1em}}l}
\tlhead{worker A} & \tlhead{worker B} \\
(1)-(3); $\fo$() on S & steal $\vt$ from A \\
& fresh stack T; $\ft$() on T \\
(4) pop() returns \textbf{null} \\
\ind restore(A.k) \\
\ind notify: $\vz$.in = 1 \\
\ind find other work ... \\
& $\ft$() returns \\
& restore(B.k) \\
& notify: $\vz$.in = 0 \\
& schedule, pop $\vz$ \\
& restore($\vz$.k): the save \\
& at (1) returns, on S
\end{timeline}
\vspace{4pt}
\hrule
\caption{The two cases: $\vt$ not stolen (top), and $\vt$ stolen by
  worker B, which finishes last (bottom)}
\label{fig:cases}
\end{figure}

Figure~\ref{fig:cases} traces both cases of \lstinline|fork2join|.

\paragraph{Not stolen}
%
One worker, A, does everything.
%
$\fo$ and $\ft$ both run on the parent's task stack, S, as ordinary
calls, and the only overhead is three round trips through the
scheduler.

\paragraph{Stolen}
%
Worker B steals $\vt$ from A while A runs $\fo$, gives $\vt$ a fresh
stack, T, and runs $\ft$ there.
%
When the pop at step 4 fails, A abandons S, and its scheduler notifies
$\vz$ that $\vo$ has finished, which brings the in-counter of $\vz$ to 1.
%
A then looks for other work.
%
When $\ft$ returns, B's scheduler notifies $\vz$ that $\vt$ has
finished, which brings the in-counter to 0.
%
So B schedules $\vz$, pops it, and restores $\vz$\lstinline|.k|.
%
The call to \lstinline|save| at step 1 returns a second time, on stack
S as before, but now on worker B, and \lstinline|fork2join| returns
there.
%
Had A finished last, A would have resumed the parent itself.

So a call to \lstinline|fork2join| can begin on one thread and return on
another.
%
The stack stays the same throughout; only the thread that runs it
changes.

\section{Why it is safe}
\label{sec:safety}

The correctness of the algorithm rests on four facts.

\paragraph{Only vertices that have never run are stolen}
%
\lstinline|fork2join| schedules $\vo$ (step 2) and $\vt$ (step 5)
before it saves their contexts.
%
At that moment, each is the vertex scheduled most recently, so it
cannot be stolen, and its own worker pops it only after the
\lstinline|swap| to the scheduler has saved its context.
%
The only vertices that thieves ever see are those like $\vt$ at step 2,
whose action is \lstinline|Init|.

\paragraph{Completion is reported from the scheduler}
%
As soon as the in-counter of $\vz$ reaches zero, another worker may
resume $\vz$ on stack S.
%
So the decrement on behalf of $\vo$ must happen after A has left S,
which is why step 4 switches to the scheduler instead of decrementing
in place.
%
Similarly, $\vo$ and $\vt$ live in the frame of $\vz$ on S, so
\lstinline|notify| reads the successor of a vertex before decrementing,
and does not touch the vertex afterward.

\paragraph{The pop at step 4 is correct even if $\fo$ moved}
%
$\fo$ returns on a different worker only if a vertex forked inside $\fo$
was stolen.
%
Thieves take the oldest ready vertex, and $\vt$ is older than any
vertex forked by $\fo$.
%
So by then $\vt$ has been stolen as well, and the pop correctly fails.

\paragraph{A fresh stack is freed only after its last use}
%
A stolen $\vt$ owns its fresh stack, and $\vt$ lives in the frame of
$\vz$, so the stack is freed when the parent's call to
\lstinline|fork2join| returns.
%
By then, $\ft$ has finished, and its worker has switched back to its
scheduler.

\section{Costs and consequences}
\label{sec:costs}

\paragraph{Costs}
%
A fork whose second branch is not stolen makes three round trips
through the scheduler (steps 3, 5, and 6), each a save and a restore,
and allocates nothing, because $\vo$ and $\vt$ live in the frame of
\lstinline|fork2join|.
%
In this case, no other worker can reach the in-counter of $\vz$, so
TaskPaRTS updates it with plain loads and stores, rather than with
atomic operations.
%
The round trips are not essential to the technique.
%
They keep the scheduler loop the only place where vertices start,
resume, and finish, and they could be optimized away.
%
A steal allocates one fresh stack, 512\,KB in TaskPaRTS, which is
touched only as deep as $\ft$ recurses, and is freed when the parent
returns from \lstinline|fork2join|.

\paragraph{Consequences for client code}
%
Code that follows a call to \lstinline|fork2join| may run on a
different thread from the code that precedes it.
%
So client code must not hold a lock across \lstinline|fork2join|, nor
expect a thread-local variable or the thread id to be the same
afterward.
%
More subtly, compilers assume that a function never changes threads,
and so may reuse a thread-local address computed before the call.
%
Then even a fresh read of a thread-local variable after
\lstinline|fork2join|, in the same function, can return the old
thread's copy.
%
Reading thread-local state through a function that the compiler can
neither inline nor treat as pure avoids this problem.
%
Other consequences follow from the use of task stacks.
%
Task stacks have a fixed size, and no guard page by default.
%
Exceptions must not escape $\fo$ or $\ft$, because unwinding cannot
cross from a task stack into the scheduler.
%
Backtraces end at the bottom of a task stack.
%
And sanitizers, such as AddressSanitizer and ThreadSanitizer, are not
told about the stack switches, so they may report false errors.

\paragraph{Implementations}
%
TaskPaRTS implements the algorithm of Figures~\ref{fig:scheduler}
and~\ref{fig:fork2join}, and adds lock-free
deques~\cite{arora98,chasele05}, elastic scheduling, instrumentation,
and CPU pinning around it.
%
Micro taskparts implements the same algorithm and nothing else, in
about 440 lines of C++, including the assembly for x86-64 and AArch64.
%
Appendix~\ref{sec:map} maps the pseudocode to both.

\bibliographystyle{abbrv}
\bibliography{main}

\clearpage
\appendix

\section{Machine contexts}
\label{sec:contexts}

A context holds the registers that the calling convention requires a
callee to preserve~\cite{sysvabi,aapcs64}, plus the stack pointer and
the address at which to resume:
%
\begin{center}
\small
\begin{tabular}{@{}lll@{}}
\toprule
 & x86-64 & AArch64 \\
\midrule
callee-saved & \texttt{rbx}, \texttt{rbp}, \texttt{r12}--\texttt{r15} & \texttt{x19}--\texttt{x29}, \texttt{d8}--\texttt{d15} \\
stack pointer & \texttt{rsp} & \texttt{sp} \\
resume address & return address & \texttt{lr} (\texttt{x30}) \\
total & 8 words & 21 words \\
\bottomrule
\end{tabular}
\end{center}
%
Caller-saved registers need no saving: to the compiler,
\lstinline|save| is an ordinary call, so the compiler has already
spilled any live caller-saved registers before calling it.
%
Figure~\ref{fig:x64} shows both functions for x86-64.
%
\lstinline|save| records, as the stack pointer to restore, the value
the stack pointer will have once \lstinline|save| has returned, and, as
the resume address, its own return address.
%
\lstinline|restore| reloads the registers, switches stacks, and jumps
to the resume address with 1 in the return register.
%
The AArch64 versions are analogous, and resume with
\lstinline|ret| to the restored \lstinline|lr|.

\begin{figure}
\hrule
\begin{lstlisting}[language={},mathescape=false]
save:                     restore:
  movq %rbx, 0(%rdi)        movq 0(%rdi), %rbx
  movq %rbp, 8(%rdi)        movq 8(%rdi), %rbp
  movq %r12, 16(%rdi)       movq 16(%rdi), %r12
  movq %r13, 24(%rdi)       movq 24(%rdi), %r13
  movq %r14, 32(%rdi)       movq 32(%rdi), %r14
  movq %r15, 40(%rdi)       movq 40(%rdi), %r15
  leaq 8(%rsp), %rdx        movq 48(%rdi), %rsp
  movq %rdx, 48(%rdi)       movl $1, %eax
  movq (%rsp), %rdx         jmpq *56(%rdi)
  movq %rdx, 56(%rdi)
  xorl %eax, %eax
  ret
\end{lstlisting}
\hrule
\caption{Saving and restoring a context on x86-64}
\label{fig:x64}
\end{figure}

\paragraph{Declaring save to the compiler}
%
Because \lstinline|save| can return twice, it must be declared as such
to the compiler (in GCC and Clang, with the attribute
\lstinline|returns_twice|).
%
Otherwise, the compiler may keep a value that changes between the two
returns in a callee-saved register, or reuse its stack slot for another
value, and the second return then sees a stale or clobbered value.

\paragraph{Fresh contexts}
%
To start a function $f$ on a fresh stack $s$, the scheduler makes a
context with \lstinline|new_context($s$, $f$)|.
%
The stack pointer of this context is the top of $s$, aligned to 16
bytes, its resume address is $f$, and its other registers are zero.
%
On x86-64, the stack pointer is a further 8 bytes down, as if a
\lstinline|call| had pushed a return address.
%
Restoring the context then starts $f$ on $s$, as if $f$ had been called
there.
%
TaskPaRTS builds fresh contexts differently, with the same effect: the
scheduler loop calls \lstinline|save|, and then overwrites the saved
stack pointer with the top of a fresh stack.
%
Restoring that context returns from that call to \lstinline|save| on
the fresh stack, where the code that follows calls \lstinline|enter|.

\section{Ready vertices}
\label{sec:deque}

\begin{figure}
\hrule
\begin{lstlisting}
type worker = {k : context, cur : vertex*,
  next : vertex*, dq : deque of vertex*}

schedule(v : vertex*) -> void {
  if (me.next != null) push_bottom(me.dq, me.next)
  me.next = v }

pop() -> vertex* {
  v = me.next
  if (v != null) { me.next = null; return v }
  return pop_bottom(me.dq) }

steal() -> vertex* {
  return pop_top(random_other_worker().dq) }
\end{lstlisting}
\hrule
\caption{Ready vertices: a private slot over a deque}
\label{fig:deque}
\end{figure}

Each worker keeps its ready vertices in a private slot, \lstinline|next|,
in front of a work-stealing deque, \lstinline|dq|
(Figure~\ref{fig:deque}).
%
\lstinline|schedule($v$)| puts $v$ in the slot, and moves the slot's
previous occupant to the bottom of the deque.
%
\lstinline|pop| takes from the slot first, and then from the bottom of
the deque, and thieves take from the top of the deque.
%
So the vertex scheduled most recently is always in the slot, where no
thief can reach it, which is the property that
Section~\ref{sec:safety} relies on.

At step 2 of \lstinline|fork2join|, the slot is empty: the current
vertex was itself popped from it, and any nested \lstinline|fork2join|
calls have taken back what they scheduled.
%
So \lstinline|schedule($\vt$)| puts $\vt$ in the slot, and
\lstinline|schedule($\vo$)| moves $\vt$ to the deque.
%
The slot is also empty whenever \lstinline|notify| or a thief schedules
a vertex, so a vertex that reaches the deque is always a $\vt$ displaced
at step 2.

Any work-stealing deque will do.
%
Micro taskparts uses a deque protected by a lock, whereas TaskPaRTS
uses a lock-free deque, chosen at compile time among the
deque of Arora et al.~\cite{arora98}, that of Chase and
Lev~\cite{chasele05}, and a deque by Yao, Westrick, Rainey, and Acar,
which supports elastic scheduling.
% TODO: cite the Yao, Westrick, Rainey, and Acar deque.

\section{From pseudocode to code}
\label{sec:map}

\begin{table*}
\small
\begin{tabular}{@{}lll@{}}
\toprule
pseudocode & micro taskparts & TaskPaRTS \\
\midrule
\lstinline|save|, \lstinline|restore| & \lstinline|mtp_context_save|, \lstinline|mtp_context_restore| & \lstinline|context_save|, \lstinline|context_restore| \\
\lstinline|vertex| & \lstinline|vertex| & \lstinline|native_fork_join_vertex| \\
\lstinline|action| & \lstinline|action| & \lstinline|continuation_action| \\
\lstinline|next|, \lstinline|dq| & \lstinline|next|, \lstinline|deque| & \lstinline|bottom|, \lstinline|deque| \\
\lstinline|loop| & \lstinline|worker_loop| & \lstinline|loop| \\
\lstinline|new_context| & \lstinline|new_context| & \lstinline|new_continuation| \\
\lstinline|enter| & \lstinline|vertex_entry| & \lstinline|enter| \\
\lstinline|notify| & \lstinline|notify| & \lstinline|parallel_notify| \\
\lstinline|fork2join| & \lstinline|fork2join| & \lstinline|fork2join| \\
\bottomrule
\end{tabular}
\caption{Names in the pseudocode, in micro taskparts, and in TaskPaRTS}
\label{tbl:map}
\end{table*}

Table~\ref{tbl:map} maps the names in the pseudocode to their
counterparts in micro taskparts and in TaskPaRTS.
%
Both are in the TaskPaRTS repository~\cite{taskparts}: micro taskparts
in the directory \path{doc/native-fork-join/microtaskparts}, and
TaskPaRTS in the files \path{src/taskparts.cpp} and
\path{include/taskparts/taskparts_internal.hpp}.
%
Both implementations number the steps of \lstinline|fork2join| as in
Figure~\ref{fig:fork2join}.

Beyond micro taskparts, TaskPaRTS adds the following.
%
\begin{itemize}[leftmargin=*]
\item Lock-free deques (Appendix~\ref{sec:deque}).
\item Elastic scheduling: workers that fail to find work suspend on a
  semaphore, instead of spinning, and are woken when there is surplus
  work.
\item Non-atomic in-counter updates when the second branch is not
  stolen (\lstinline|clone_fast|).
\item A persistent pool of workers, and a sink vertex that marks the
  end of each top-level computation, in place of the
  \lstinline|done| flag.
\item Statistics, logging in a format that the Perfetto trace viewer
  reads, and CPU pinning with hwloc.
\end{itemize}

\end{document}
